\documentclass{article}
% >>> COMPILE WITH XeLaTeX (Overleaf: Menu -> Settings -> Compiler -> XeLaTeX). <<<
\usepackage{graphicx} % Required for inserting images
\usepackage{xcolor} % Must be loaded before soul
\usepackage{soul} % For highlighting text
% inputenc/fontenc intentionally omitted: XeTeX is natively UTF-8 and fontspec provides the encoding.
\usepackage{amssymb} % \therefore used in the metaphysics premise blocks
\usepackage{newunicodechar}
\usepackage{tikz}
\usetikzlibrary{positioning}
\usepackage{fontspec}           % For Greek text support with XeLaTeX/LuaLaTeX
\usepackage{polyglossia}        % Multilingual support
\setmainlanguage{english}
\setotherlanguage{greek}

\usepackage{geometry}
\geometry{margin=1in}

\usepackage{csquotes}           % Block quotes
\usepackage{enumitem}           % Better lists
\usetikzlibrary{arrows.meta, positioning, shapes.geometric}

\usepackage{fancyvrb}           % Verbatim environments for ASCII diagrams
\usepackage{setspace}           % Line spacing

% --- Greek support (XeLaTeX) ---
\newfontfamily\greekfont{Gentium Plus} % good polytonic Greek; try Times New Roman if needed
\newcommand{\gk}[1]{{\greekfont #1}}

% RODA chain glyphs: subscripts A_0..A_4, plus the arrow / logic / therefore symbols used in the text
\newunicodechar{₀}{\ensuremath{_0}}
\newunicodechar{₁}{\ensuremath{_1}}
\newunicodechar{₂}{\ensuremath{_2}}
\newunicodechar{₃}{\ensuremath{_3}}
\newunicodechar{₄}{\ensuremath{_4}}
\newunicodechar{→}{\ensuremath{\rightarrow}}
\newunicodechar{∧}{\ensuremath{\land}}
\newunicodechar{∴}{\ensuremath{\therefore}}
\newunicodechar{′}{\ensuremath{{}^{\prime}}} % prime, e.g. A_0' (next-cycle deposit)

\usepackage[backend=biber,style=authoryear]{biblatex}
\addbibresource{references.bib} %Add References page for in text citation
\usepackage[colorlinks=true, linkcolor=blue, urlcolor=blue, citecolor=blue]{hyperref} %enable hyperlinks to be colored and interactable
\usepackage{changepage} %create individual blocks of text that can be altered apart from general document.

\usepackage{comment}

\newcommand{\ra}{\ensuremath{\rightarrow}\ }

\newcommand{\inlinenote}[1]{\par\noindent
  \fcolorbox{orange}{orange!20}{\parbox{0.97\linewidth}{#1}}\par\noindent}

\title{Part III --- The Stalled Chain: Boredom, Engagement, and the Design of Virtual Learning Environments \\ \large DRAFT v1 (FCDP full run)}
\author{}
\date{2026-07-07}

\begin{document}
\maketitle

\section{Introduction: Attention, the Digital Habitat, and the Reanalysis Stance}

Where the account so far has moved from the inside outward---from the motion of the soul to the designed worlds that solicit it---the present Part moves back: from the frame to the laboratory, from the mechanism to the measurements, from the claim that a designed world runs or stalls the actualization of desire to three studies and one living course in which participants were instrumented, students were asked, and a designed world did exactly that. The experiments, moreover, are the author's own: the virtual learning environment under study was built for an undergraduate biomedical engineering course; the studies that instrumented it---one physiological, one oculometric, one phenomenological---were conducted and published with its builder's hand in each of them; and the course itself, run wholly online and at the scale of a full enrollment, has now yielded three terms of student testimony about life inside the world it delivers. What follows is therefore not a new experiment but a reanalysis of the whole record: the published findings and the primary data are read again, through the frame this dissertation has built, so that the data yield the conclusions while the frame supplies what data alone can never supply---the how and the why of their arising.

Three goals govern the Part, and they are best named at the threshold. The first is boredom: to analyze the student's boredom phenomenologically and rhetorical-ontologically, so that what the instruments recorded as the rising and falling of signal becomes legible as the structure of an attunement. The second is engagement: to assess what the surveys and the studies called engagement against the dissertation's own methodology, and to say, with the frame's precision rather than the instrument's, what that word was measuring in each of its appearances. The third is design: to show that a frame able to explain why virtual worlds bore their students can direct the making of worlds that do not, so that diagnosis, carried through to its end, becomes a discipline of design. Boredom explained, engagement assessed, design instructed---the three goals are one movement performed three times, each time at a greater depth of the same chain.

One fact about the corpus must stand at the door of the Part, because everything that follows passes through it. Between the pilot studies and the deployment at scale, the course moved: virtual reality remained available to any student who sought it, but a full-sized class, run wholly online and asynchronously, could not be scaled into headsets, and so the environment met nearly all of its students as a three-dimensional world presented on a two-dimensional screen. Far from a fall away from the true experiment, the move was the discovery of the actual one---the discovery of which conditions a designed world must actually satisfy when it leaves the laboratory for the enrollment list.

Why should engagement with digital artifacts be analyzed at all? Because the question deserves its own bluntness, the answer should begin where it in fact begins, a century before the screen. William James, whom Gloria Mark sets at the head of her account of the modern attention span, held that attending is not a lamp the mind plays over a finished world but the very act by which a life takes its shape: ``Millions of items of the outward order are present to my senses which never properly enter into my experience. Why? Because they have no interest for me. My experience is what I agree to attend to. Only those items which I notice shape my mind---without selective interest, experience is an utter chaos.''\footnote{William James, qtd. in Gloria Mark, \textit{Attention Span} (Toronto: Hanover Square Press, 2023), 30.} Experience, on this account, is not what happens to a person but what is attended to while it happens---a sentence whose weight falls, and is meant to fall, on the verb. Mark carries the claim into the present tense of every pocket: walking through a garden with a phone in hand, she observes, it is the exchange of messages that survives as the record of the walk---``It is the texts that have entered into the record of my experience, and not the softness of the ground, the trill of the warbler or the scarlet of the azaleas.''\footnote{Mark, \textit{Attention Span}, 30.}---while the garden, present to the senses and unattended by the mind, never properly enters the experience at all. What holds the attention holds, in the end, the life.

And attention is now held by screens---held briefly, held often, and held in numbers that permit a census of the holding. Mark's own logging studies, replicated across research teams and workplaces, find that ``in the last several years, every day and all day in the workplace, people switch their attention on computer screens about every forty-seven seconds on average.''\footnote{Mark, \textit{Attention Span}, 93--94.}, and by 2016 the median stretch of attention on any screen had fallen to forty seconds.\footnote{Mark, \textit{Attention Span}, 94.} Whatever else the digital artifact may be---tool, medium, world---it is the place where attending now factually lives, in tenancies measured in seconds and ended by a switch, several hundred times in a working day. To study human engagement with digital artifacts is not, then, some regional curiosity of educational technology but the study of attention in its native habitat, conducted at the very site where the constitution of experience is presently decided.

An objection stirs at this point, and it should be given its full strength: the Jamesian warrant proves too much, for by its light every medium---book, lecture, window---would equally constitute experience, and nothing particular would follow for the screen. What answers the objection is the concentration. The claim is not that screens alone shape what enters a life; the claim is that screens are where the shaping now happens with measured regularity and measured brevity, so that the forty-seven seconds function not as a metaphysical thesis but as a census return. Where attention lives is an empirical question, and it has received an empirical answer.

The dissertation's frame, it should be said, speaks the same conclusion in its own tongue. What is attended is what enters the chain---what becomes \textit{phantasma}, takes on \textit{doxa}, and issues in act---and what enters the chain, entering repeatedly, sediments into the \textit{hexis} that a student carries out of a course and into a practice. Attention, so placed, is not the antechamber of learning but learning's first act, and a world that cannot hold attention cannot form anyone at all.

How attention is studied has meanwhile turned on its own axis, and the turn supplies this Part with its instruments. Where an earlier psychology of ability asked its subjects and tested them, the present science senses them: ``surveys and tests are no longer needed when sensors can unobtrusively detect people's physiological signals and tracking can detect digital traces of their online behaviors.''\footnote{Mark, \textit{Attention Span}, 154.} Mark reports the turn as an accomplished fact of the field, and the corpus of this Part had enacted it in miniature before the frame ever arrived: across the three studies the heart-rate measures were dropped as uninformative, the electroencephalogram---too noisy in laboratory practice---was demoted to corroboration, the eye tracked at a hundred and twenty hertz was elevated to the primary instrument, and the third study turned at last to the phenomenological interview. On its own methodological legs, that is, the program walked from asking toward sensing, and from sensing toward the description of lived experience---a walk driven not by theory but by what the instruments could and could not see, and what they could not see singly is precisely what this Part will give its name.

The same instrumentation carries a mirror, and the mirror sets the stakes. Industrially deployed, the sensors that might diagnose a stalled attention are used instead to seize it at its weakest: ``If you receive targeted notifications going after your lower-order emotions when you are low in attentional resources or in a bored or a rote state, when you are more susceptible to distraction, you can't help but respond to them.''\footnote{Mark, \textit{Attention Span}, 159--60 [****** page range unverified].}. The notification engineered for the bored is the dark twin of the lesson designed for the engaged, since both proceed from knowledge of the stalled chain---the one to exploit the stall, the other to repair it---and the symmetry of the knowledge is the strongest reason for a pedagogy to hold it. A pedagogy that declines to understand why elements bore does not thereby keep its hands clean; it merely leaves the understanding, and the students, to those who monetize both. The mirror will return at the Part's end, where the norms of design are weighed; here it need only set the stakes.

A second objection cuts closer, because it cuts at the corpus itself: if surveys and tests are superseded by sensors, what standing have the three terms of survey testimony on which the deployment analysis will lean? The concession is real, and it is bounded. No single channel---sensed or asked---ever suffices for the phenomenon at issue here, and the frame's central methodological rule, stated in the next section and worked throughout, is precisely that the subtlest of boredom's forms shows itself only where two channels disagree. The survey is one voice among the channels, indispensable and insufficient; the instruments are another; and the analysis that follows listens, above all, for their divergence.

One misreading, finally, must be blocked before the corpus is opened, because the very vocabulary of the frame invites it. To say that a designed environment runs or stalls the actualization of desire sounds, to a certain ear, like the oldest behaviorism in a new costume---``To Skinner, human cognition is a fallacy. He argued that our environments condition us in how we behave in our daily lives.''\footnote{Mark, \textit{Attention Span}, 252--53.}. The resemblance does not survive the frame's own grammar. In the account advanced here, the designer and the participant divide the causes of the act as the producing art divides them from the using art: the designer furnishes the matter and the form of the encounter---the rooms, the images, the affordances---and by furnishing them occasions the act, while the efficient and the final cause remain the participant's own, or else there is no act at all but only motion. A world conditions nothing; a world offers. What the participant does with the offer---whether the \textit{phantasma} is taken up, whether \textit{doxa} commits itself, whether anything is done---is the participant's actualization and never the environment's output, so that cognition, far from being a fallacy, is the very site where the offer is accepted or refused. When a later section says of the deployment that its world gave students nothing to do, the grammar of the complaint should therefore be heard exactly: not that the world failed to produce behavior in its occupants, but that it withheld the conditions under which a student could complete an act of one's own. The designer's power is real and it is bounded---real, because the offer shapes everything that can happen; bounded, because the offer is all that it is. Design, on these terms, is rhetoric, and rhetoric persuades agents; it does not program mechanisms.

The corpus, then, stands at two poles. At the first stand the laboratory and the pilot: between 2022 and 2024, three published studies instrumented the environment's early life---a physiological pilot that watched three participants' brains and hearts across three seventeen-minute films, an eye-tracking study that followed twelve participants' gazes at a hundred and twenty hertz, and a phenomenological evaluation that asked and interviewed a pilot cohort while four formats of the same clinical footage competed for their preference. Around the studies stands the raw record itself, newly assembled for this Part: eight participants' electroencephalograms, their headset telemetry and session recordings, and---most valuable of all, as it will emerge---their written answers, film by film, to seven blunt questions, among them the question of how long each film felt. At the second pole stands the deployment: three terms of the live course, two hundred forty-one students, eight hundred seventy-three written answers, and a world reached through a downloadable application on the student's own computer---or, for nearly half of every term with remarkable constancy, not reached at all, in favor of the plain videos. Between the poles falls the move already set at the door: presence spoke for the headset, feasibility spoke for the screen, and the course, like most of its students, followed feasibility.

Read cold, the record already leans toward the frame, for its instruments disagreed, and the disagreement is the door. On the same clinical film, the earlier study's electroencephalogram read a profile that resembled boredom while the later study's gaze-variance read engagement with statistical confidence, and the participants meanwhile called themselves interested while the felt clock stretched, the clean recording's owner feeling seventeen minutes as twenty-five to thirty. No bipolar scale can hold that record together: each instrument, taken alone, tells a different story, and the program, to its credit, kept the contradiction rather than smoothing it away. The frame will read the contradiction in one move---a surface occupied, a depth left empty---and that reading will organize everything which follows: the laboratory's forms of boredom, the deployment's flights and frictions, and the design levers that answer them both.

The route through the Part runs so: first the method of reanalysis, stated together with its limits; then the laboratory, where the forms of boredom acquire their empirical faces; then the deployment, where a designed world met two hundred forty-one students and, for the most part, showed them things; then the synthesis, where the corpus is read as one record and the question of why elements bore receives its answer; and last the return to design, where the lesson that to design a world is to design the conditions under which a chain runs or stalls is brought home to the classroom with the corpus's own evidence in hand. The data are already in hand. What they awaited was the frame.

\section{The Method of Reanalysis}

The method of this Part is the method of the dissertation carried to data, and its first word is a disclaimer of novelty: nothing in the frame is revised here, and nothing is added to it except an application layer---a set of stated jurisdictions by which each engine of the account meets each kind of evidence on terms fixed in advance. Four engines do the work. The analysis of boredom runs on the apparatus drawn from Heidegger's lecture course: the three forms of the attunement, their two structural moments, and a diagnostic grid that asks its questions in a fixed order---whether a determinate culprit can be named; whether an occupied surface stands over an emptied depth; whether the passing of time has ceased altogether while the while stretches. The analysis of design runs on the lens drawn from phenomenological design theory, which reads an environment as a set of conditions that afford or foreclose the chain's stages, and which will be held to its predictions when the deployment's complaints are counted. The game-analysis method developed for Part II's cases serves here in a deliberately bounded role, since its protocol wants a recorded episode of play---a trace in which an actualization can be followed beat by beat---and only the laboratory recordings supply one; where the data are testimony rather than trace, its vocabulary serves and its worksheet stays shut. Beneath all three, the metaphysics and the methodology of the earlier Parts govern the words: the chain from horizon to act, the union with its two poles of world-disclosedness and co-disclosedness, the economy of the image named by \textit{phantasia}-load. Jurisdiction, not doctrine, is what this section settles---which engine may speak of which data, and in what register of confidence.

That register is tiered, and the tiers are worn openly in the prose. Findings of the published studies are asserted as findings; the raw archive of the author's own laboratory work, which has not been published, is asserted with that status marked; and every reading that carries an empirical result over into the frame's categories is flagged as the frame's reading---a projection offered for its explanatory yield, never smuggled in as one more datum. The reader will therefore meet three voices in what follows: the record speaking, the record's unpublished margin speaking, and the frame speaking about both. Keeping them distinct is not decoration. It is what entitles the Part to its conclusions.

The founding constraint on all four engines comes from the lecture course itself, and it must be stated before any instrument is trusted, because it disciplines them all. Heidegger opens the analysis of boredom by refusing the very posture an empirical program most naturally assumes: ``We shall not speak at all of `ascertaining' a fundamental attunement in our philosophizing, but of awakening it.''\footnote{Martin Heidegger, \textit{The Fundamental Concepts of Metaphysics: World, Finitude, Solitude}, trans. William McNeill and Nicholas Walker (Bloomington: Indiana University Press, 1995), 60 / GA 29/30: 90.} An attunement is not an object lying present for inspection; it is the way a situation already has us, and the attempt to seize it as an object alters or destroys the thing seized---``any so-called objective ascertaining of a fundamental attunement is a dubious, indeed impossible undertaking.''\footnote{Heidegger, \textit{Fundamental Concepts}, 60 / GA 29/30: 90.} Taken seriously, the constraint does not forbid measurement; it relocates what measurement measures. Electrodes, eye-trackers, and questionnaires never read the attunement itself; they read comportment---the fidgets and glances and answers in which an attunement casts its shadow---and the analysis that follows treats every signal accordingly, as shadow-evidence whose bearer must be reconstructed, never observed. A second consequence follows for classification. If boredom comes in forms rather than in quantities, then the forms are told apart by their structure and not by their intensity: the depth of an episode, the lecture course insists, ``can be read off only from the specific constitution of boredom itself in each case.''\footnote{Heidegger, \textit{Fundamental Concepts}, 130 / GA 29/30: 195--97.} A scale that runs from engaged to bored therefore begs the decisive question in advance, for it presumes that the phenomenon has one dimension, when the second form's whole character is to hold two dimensions apart---a surface that remains occupied and a depth that has been left empty. Hence the rule on which this Part's deepest claims will stand or fall: the second form of boredom is detectable only as the disagreement of a surface channel and a depth channel, never from either alone, and never from an instrument built to collapse the two into one score. The rule carries a corollary for testimony that must be honored even when it is inconvenient: the subject of the second form sincerely denies being bored, reading the occupied surface as the whole story---so self-report is a channel to be listened to, and never a verdict to be accepted.

A caution of category completes the constraint, and it is constitutive of the method rather than an apology appended to it. A seventeen-minute video does not disclose beings as a whole; a course module does not visit the profound boredom in which a world hangs suspended; and nothing in this Part pretends otherwise. What transfers from the lecture course to the laboratory is the structural grammar of the forms---culprit-structure, emptiness-structure, temporal structure---taken as a heuristic that must earn its keep case by case. The first two forms transfer cleanly, since their grammar is episodic on its face. The third transfers only as an analogy, explicitly marked wherever it appears, for a collapsed gaze may resemble the stance of profound boredom without any claim that a participant underwent a metaphysical disclosure between one film and the next.

With the constraint and the caution in place, the corpus's channels can be sorted in advance, so that no instrument acquires authority by accident. As surface channels the Part admits the gaze metrics, the recorded comportment, and every self-report of engagement or presence, closed or open---all of them readings of the occupied surface. As proxies for depth it admits three, unequal in strength: the electroencephalographic profile, which enters with corroborating weight only; the felt-duration reports, which enter as the strongest depth evidence the corpus owns, since the stretching of the while is the one signature that the second form cannot suppress; and, weaker still, the direct self-ratings of boredom and of hollowness---the closed rating where an instrument asks about boredom separately from engagement, and the hollowness that leaks through open testimony---admitted as depth traces only where they stand against a surface-positive answer from the same voice, and always under the denial-caution stated above, since the second form's subject reads the occupied surface as the whole story. Divergence---the method's true object---is admitted wherever a surface and a depth speak of the same episode: instrument against instrument in the laboratory, item against item within a single laboratory questionnaire, and, in the deployment testimony, only within a single student's answers, where praise and hollowness share one sheet. And some things are admitted as nothing at all: the content-quiz scores, which sat at ceiling and so carry no signal; the platform's automated item statistics, which assign correct answers to survey questions and are artifacts of the export rather than psychometrics; and any composite that would average engagement and boredom into a single number, which the rule above disqualifies in principle. The sorting is itself falsifiable, and should be seen to be: nothing guarantees that surface and depth will ever disagree, and had every channel of this corpus agreed with every other, the second form would simply have found no purchase here. The frame's design lens, likewise, is made to stake predictions---which frictions a student population would name---that the deployment section counts against the record. A frame that could not have failed would explain nothing by succeeding.

The deployment testimony, finally, was prepared for this method by a coding protocol whose reliability is stated rather than assumed. The instrument is a fourteen-item course questionnaire, stable across the three terms studied. Three content items sat at ceiling and were excluded, and one item is a survey preamble carrying no data. Two closed items carry the experience data: whether the student was able to use the platform, and whether it produced a feeling of immersion or presence, with a video-only answer available. Four open items invited prose; four demographic items close the form. Two hundred forty-one students answered across the terms, and their eight hundred seventy-three open responses form the coded corpus. The codebook holds eight families: boredom-form signatures, with sub-tags for the named frictions; flight patterns; temporal texture; design conditions keyed to the chain's stages; involvement-channel vocabulary; incorporation-positive signatures; suggestions; and an open family for emergent codes. Every response was coded at least once. A subsample of one hundred forty-three responses was coded twice, blind, and the two passes agreed at a mean Jaccard of .878 over full code sets; the four disagreements that survived rule-governed merging were adjudicated by hand and logged. Names were replaced by codes throughout, and references to instructors were masked. The record's limits are equally part of the record. The survey is one self-report channel, so the second form can appear in it only as an echo---a divergence-shaped pattern inside a single voice---and never as a detection. The questionnaire is retrospective, so the temporal texture of episodes is largely mute in it, and the laboratory archive must carry that thread. No outcome measure survives the ceiling, so nothing here speaks of learning gains. And the terms differ in composition, so their comparison stays directional. Within those limits, the corpus is ready; what remains is to read it.

\section{The Boredom Experiment Re-Read}

The record comes first, plainly stated, with its reliability flags attached. The first study was a within-participants physiological pilot with three participants. Each watched three seventeen-minute videos: a boring control (a 1989 word-processor tutorial), an interesting control (an alien-conspiracy video), and the clinical stimulus (360-degree footage of spinal-deformation surgery). An electroencephalographic cap recorded four active channels (F3, F4, P3, P4) at 200 Hz, targeting the default-mode network's alpha (8--12 Hz) and theta (4--8 Hz) bands; heart rate, gaze, and pupil diameter came from the headset's sensor suite. Its clean findings are two. All three participants felt the clinical video as longer than its seventeen minutes---the one clean recording's owner felt it as twenty-five to thirty---and self-reported engagement ran high, with two participants rating interest at eight of eight. Its headline physiological finding is suggestive only: the clinical video's cortical profile more closely resembled the boring control than the interesting one, but only one of the three recordings was clean, one carrying movement artifacts and one impedance problems. The heart measures were inconclusive, and learning itself, despite the study's title, was never measured. The second study replaced the brain with the eye. Twelve participants watched the same three videos while gaze was tracked at 120 Hz; the analysis took pupil dilation, gaze position, deviation from center, and---its strongest measure---the variance of gaze deviation under a five-point sliding-window median, with analysis of variance at $\alpha = 0.05$. A four-person webcam sub-study validated the approach against consumer hardware (average cross-correlation 0.874). The presentation order was changed to interesting, then clinical, then boring, because a boring-first order had left participants too exhausted for what followed---a fact that will matter later. The findings: gaze-variance on the clinical video was statistically distinct from the boring video ($p = 0.0004$) and resembled the interesting one, while every other feature-by-video comparison stayed nonsignificant; pupil dilation exceeded the boring baseline for the interesting video in six of twelve participants; gaze-variance ran lower for interesting than boring in eight of twelve; and the study drew its own distinction between two boredoms---a high-arousal boredom whose gaze roves the screen searching for something to attend to, and a low-arousal boredom, the zombie stare, in which variance collapses, observed in two participants. Around the published studies stands the raw archive of the same program, reported here with its unpublished status marked and its participants anonymized. Its instrument is a seven-item questionnaire answered after each video: fatigue before; boredom and engagement, each on a nine-point scale; minutes until boredom began; a sleep-fight rating; felt duration; fatigue after. Eight participants completed it for all three videos. The aggregates: mean boredom of 7.0 for the boring video, 6.0 for the clinical, 4.5 for the interesting; felt duration of twenty minutes or more, against seventeen actual, in five of eight for the boring video, six of eight for the clinical, five of eight for the interesting---sixteen of twenty-four episodes dilated, and every felt value in this record that cleared the clock also cleared twenty minutes; boredom onset within five minutes in seven of eight for the boring video, two of them under forty seconds; and fatigue after the boring video at or above its starting level in six of eight.

Read through the frame, the boring stimulus behaves exactly as becoming-bored-by predicts, and its behavior calibrates everything else. The first form of boredom has a nameable culprit---this tutorial, this drone---and a proto-definition that the lecture course states with both structural moments in one breath: ``that which bores, which is boring, is that which holds us in limbo and yet leaves us empty.''\footnote{Martin Heidegger, \textit{The Fundamental Concepts of Metaphysics: World, Finitude, Solitude}, trans. William McNeill and Nicholas Walker (Bloomington: Indiana University Press, 1995), 87 / GA 29/30: 130--32.} Both moments have empirical faces here. Being left empty (\textit{Leergelassenheit}) shows as the collapse of interest ratings toward the floor and the near-immediacy of onset: seven of eight participants were bored within five minutes of a seventeen-minute film, two of them within forty seconds, which is to say that the tutorial's offer was refused almost at the threshold, and everything after the refusal was endurance. Being held in limbo (\textit{Hingehaltenheit}) shows in the sleep-fight column and in the stretched while---twenty and twenty-five and thirty felt minutes against the seventeen actual, one participant reaching a full felt hour. And between the two moments runs the visible counter-move that the eye-tracker caught as high variance: the searching gaze, roving the screen for anything to attend to, is \textit{Zeitvertreib} made measurable, the passing-of-time in its restless first-form costume. The lecture course had already named the mechanism whereby the counter-move fails: ``Passing the time is a driving away of boredom that drives time on.''\footnote{Heidegger, \textit{Fundamental Concepts}, 94 / GA 29/30: 141--42.} Passing the time drives time on and thereby keeps the dragging while in view; the search sustains the very emptiness it flees; and the eye that cannot find anything is, in its motion, the portrait of an offer that gives nothing to complete.

The clinical stimulus is another matter, and it is the section's center, because on it the instruments disagree. The cortical profile of the first study read the surgical footage as boring-like; the gaze of the second study read it as engaging-like with the corpus's one emphatic significance; the participants called themselves interested; and the clean recording's owner felt the seventeen minutes as twenty-five to thirty. The frame reads the disagreement as the signature of the second form, being-bored-\textit{with}, whose whole structure is an occupied surface over an emptied depth [a projected reading, and marked as such]. In the lecture course's paradigm---an agreeable dinner party after which one finds, at home, that one was bored after all---the passing-of-time no longer sits beside the activity as cigarette or doodle; the occupation itself becomes the passing-of-time: ``It is not smoking as an isolated occupation, but our entire comportment and behaviour''\footnote{Heidegger, \textit{Fundamental Concepts}, 112 / GA 29/30: 169--70.} An engaged-looking gaze over a hollow cortical hour is precisely that structure rendered in channels: the surface busy, the depth left empty, and no culprit anywhere to blame, for the second form's emptiness, as the course insists, ``what is boring does not come from outside: it arises from out of Dasein itself.''\footnote{Heidegger, \textit{Fundamental Concepts}, 128 / GA 29/30: 192--94.} The objection writes itself, and it deserves its strength: the two instruments belong to two studies, two years, two cohorts, and their disagreement might be an artifact of the seam. The concession is real and it is bounded, twice over. First, the keystone does not need the electroencephalogram: within the later study's own logic, felt duration against gaze supplies the two channels, and a gaze that reads engaged over a while that stretches past its clock is the same divergence without the seam. Second---and this is the raw archive's singular contribution---one participant staged the divergence inside a single instrument. Asked two questions instead of one, the participant coded S07 in the archive rated the clinical film at boredom six \textit{and} engagement seven, and the interesting film at boredom seven \textit{and} engagement six. Given separate voices for surface and depth, the answer sheet refused the bipolar axis on its own initiative; the two layers, permitted to speak separately, spoke differently. That is the two-channel rule confirmed from inside testimony itself, and it is why no instrument built on a single engaged-to-bored dimension could ever have seen what this record shows.

The felt-duration column deserves its own reading, because it is the corpus's cleanest depth evidence and its most Heideggerian. Sixteen of twenty-four episodes were felt as longer than their seventeen minutes, and the dilation crosses every stimulus category: the boring film stretches, the clinical film stretches, and the interesting film stretches for five of eight participants too. Boredom's temporality, the course argues, is not a matter of the clock running slowly but of the while itself lengthening for the one held in it---``In boredom, the while [\textit{Weile}] becomes long [\textit{lang}].''\footnote{Heidegger, \textit{Fundamental Concepts}, 152 / GA 29/30: 229.}---and the drag is a relation to time rather than a property of the timed: the interval oppresses us ``not because the progress of time is slow, but because it is \textit{too} slow.''\footnote{Heidegger, \textit{Fundamental Concepts}, 97 / GA 29/30: 146--47.} What the column shows, read so, is that the laboratory situation as such---strapped into a headset, held before a mandated film with nothing to do but watch---is already a holding apparatus, a small architecture of the held-in-limbo whose grip tightens or loosens by stimulus but never quite releases. The admission must then be paid for, because it presses on the keystone's fallback: if dilation alone were the depth criterion, the interesting film---engaged-reading gaze in eight of twelve, stretched for five of eight---would qualify as well. The fallback therefore stands on a conjunction, not on the stretch alone: an engaged-reading surface over a depth that both lengthens \textit{and} rates itself bored, which the clinical film satisfies at a mean of six against the interesting film's four and a half---and, at its sharpest, on the answer sheet that carried both poles at once. The extremity makes the structure vivid: one participant, bored within twenty-five seconds, felt the seventeen minutes as an hour, and for the same participant the clinical film did the opposite work, restoring fatigue from six to three---the while releasing where interest took hold. What is boring, at the limit, is not this tutorial or that surgery. The course's terminal thesis names it: ``What is boring is neither beings nor things as such [\ldots] but temporality as such''\footnote{Heidegger, \textit{Fundamental Concepts}, 158 / GA 29/30: 236.}

At the record's low-arousal pole stands the zombie stare, and with it the analysis must grow careful, because the third form of boredom enters this Part only as an analogy and is marked as one. In two participants the searching variance collapsed: the gaze stopped roving, the counter-move ceased, and the posture resembled---resembled---the stance the course reserves for profound boredom, in which all passing-of-time has become powerless and what remains is ``a being compelled to listen''\footnote{Heidegger, \textit{Fundamental Concepts}, 136 / GA 29/30: 205.} No claim follows that a participant underwent a metaphysical disclosure between films; the category caution forbids it, and the data do not need it. What the record adds instead is the exit that the analogy alone would miss. One participant fell asleep during the boring film, rated the sleep-fight at nine, and felt the seventeen minutes as five---the column's most extreme contraction. Where dilation is the signature of wakeful holding, sleep is not boredom's maximum but its escape: the entranced while requires a witness who remains held, and when the witness slips away, the long while vanishes with the sleeper. The methodological consequence is concrete, and it is double. A felt-duration probe saturates strangely at the sleep boundary, so any classifier built on these signatures needs sleep-exit as its own terminal category, distinct from the deepest boredom rather than continuous with it; and a collapsed gaze is by itself ambiguous between absorption and abandonment, so no classifier may read low variance alone---the companion channels, sleep-fight and felt time, must disambiguate what the stilled eye cannot.

Two features of the record remain, and they close the laboratory reading by turning it against the schema that most naturally organizes such experiments. The first is carryover. The second study reordered its films because boring-first had exhausted its participants, and the raw archive confirms the pattern from inside: fatigue after the boring film met or exceeded its starting level in six of eight participants, twice climbing from a three to a seven and an eight. Boredom, whatever else it is, depletes---an aftermath the lecture course, which treats the attunement as disclosive, never theorizes, and the one place in this corpus where the traffic runs from the laboratory back toward the philosophy. Design will have to respect it: pacing, sequence, and washout are boredom variables, not logistics. The second feature is the anomaly pair, and it is the record's quiet refutation of stimulus-essentialism. One participant entered the clinical film already exhausted, at fatigue seven, and rated it boredom nine and engagement one, the film felt as thirty minutes; the same participant then met the boring tutorial---the corpus's designated dullness---and rated it boredom one, engagement seven, felt eight minutes, fatigue falling from seven to one. Another participant was bored by the interesting film within a minute while giving the tutorial an engagement of five. The same footage bores, engages, restores, and repels, depending on the episode's temporal and dispositional situation; nothing in the stimulus fixes its effect. The cause-effect schema fails empirically exactly where the lecture course dismantles it conceptually, and the failure is the reanalysis's warrant: what the instruments were watching was never a property of three films, but the fortunes of a chain---solicited, held, emptied, completed, or abandoned---whose structure the next sections carry out of the laboratory and into a deployed world.

\section{The VLE in Deployment}

The pilot evaluation had ended on a tension it could not resolve from inside. Seventy percent of its cohort reported a greater felt presence in the headset formats, and forty percent of the same cohort judged the plain two-dimensional video the most beneficial way to learn---convenience and ease carrying 44.5 percent of the stated reasons, physical discomfort and disorientation another 38.9---while the library's loanable headset, the study noted, went unused by every student in the class. Those with more prior experience of virtual reality proved more likely to prefer the flat screen, and the minority who preferred the headset praised, above all, embodied first-person presence and the removal of distraction. Presence spoke for one route and preference for another, and between the pilot and the deployment the institution resolved the tension exactly as most of its students had: the course, run wholly online, asynchronously, and at full enrollment, could not be scaled into headsets, and so the environment went to its public as a three-dimensional world on the student's own screen, the headset remaining a permitted rarity. One wording must be settled at the door: the course's survey, drafted in the platform's earlier life, still calls the application the VR platform, and the students answer in kind; in the deployment terms at issue here the phrase names a desktop experience, and the analysis therefore speaks of app-users and of video-only students, letting the survey's own vocabulary appear only inside quotation. Nothing in the move was a loss to the inquiry, for a designed world that must survive the enrollment list, the download, and the student's own laptop is the design problem in its adult form, and the three terms of testimony it produced are the dissertation's richest evidence.

The record, plainly stated. Three terms were surveyed---102, 79, and 60 students; 241 in all---on one stable instrument, and their 873 open answers form the coded corpus read here. The strata barely move from term to term: about half of each class used the desktop application, and between 43 and 46 percent of every term used only the fallback videos, a constancy so flat that the term-to-term association is negligible. Among app-users, the immersion question drew yes from 64, 82, and 67 percent across the three terms---roughly the pilot's 70 percent, returned at scale and on a screen. Temporal texture is nearly mute in the corpus (one drag response in 873), as a retrospective instrument would predict; term compositions differ; and one duplicated response across terms is flagged in the record. Those are the limits. Within them, the corpus speaks with unusual clarity.

Where the world foreclosed, it foreclosed in two zones, and the zones divide the strata almost cleanly. The video-only half's complaints cluster before the world: installations that fail, sign-ins that refuse, machines that cannot carry the build---``More inclusive for macbook users :(''\footnote{Student W25-R042, Winter 2025 survey, platform item.} runs one answer, and sixty-one students report never having been able to use the platform at all---while the app-users' complaints cluster inside the world: navigation and wayfinding (25 mentions, seven of them arriving together with requests for guidance), latency and performance (27), crashes (24). Read against the design lens's pre-deployment walkthrough, the inside zone is a scorecard of confirmations---the read had predicted the buffering and the wayfinding, and the students confirmed both---while darkness and viewing distance, also predicted, surface only weakly, and the predicted occlusion of one avatar by another goes unattested for the plainest of reasons: a solo desktop session has no other avatars to occlude. But the scorecard's honest lesson lies in what the walkthrough could not see. Crashes, installations, and hardware access---24, 22, and 22 mentions---dominate the friction record, and none of them were predictions, because a developer's machine already runs the build. The threshold, it turns out, is the design's first and largest fact: availability is the zeroth condition of any designed world, the horizon before the horizon---and this one shut a quarter of the class out outright while taxing the entry of everyone else. To object that a student's laptop is not the designer's fault is to mistake the lesson. A world delivered as a download \textit{is} its installer, as much as its rooms; the approach belongs to the architecture; and a threshold that turned away a quarter of the class---and, counting the flight it occasioned, cost the world nearly half its public---has foreclosed more actualization than any dark room inside.

For the half that entered, the world's inner economy has a precise shape: everything to see, nothing to do. Interest and content dominate the positive record---159 mentions, twenty-one of them in the company of the corpus's incorporation-positive language, which accompanies talk of things to do exactly never---for the footage itself, the operating rooms and procedures, is what students consistently name as the value. Action is the missing column. Mentions of things to \textit{do} number 36, not one of them accompanied by incorporation-positive language, and 51 answers ask outright for interactivity the world does not have: one student wants it ``more interactable and had options on which route you want to take rather than just being playback of videos.''\footnote{Student W25-R056, Winter 2025 survey, platform item.} The passivity has its own testimony---``there was a lot of space of inactivity. Videos were unnecessarily long.''\footnote{Student W25-R025, Winter 2025 survey, improvements item.}---and its sharpest formulation names not the length but the missing object: ``Some of them are many hours long, and I'm not sure what exactly to be looking for during that time.''\footnote{Student S25-R035, Spring 2025 survey, improvements item.} There is the withheld \textit{orekton} in a student's own words: hours of material and nothing determinate to want from it, attention solicited by the image and then left standing, with no task in which the looking could complete itself. In the chain's terms the diagnosis writes itself: the world reliably produces the charged image and the committed interest, and reliably withholds the act; it runs the chain to the threshold of completion and stops; and a world that can only be watched is a world in which the watching must eventually confess that it has nothing further to become.

The confession, when it comes, takes the form of flight, and the flight has two anatomies that the closed answers separate cleanly. Sixty-one video-only students never got in: theirs is an excluded flight, decided by the threshold. But forty-seven students answered that they \textit{could} use the platform and stayed with the videos anyway, and their prose does explicit arithmetic: ``It's not worth the effort to even learn how to use it and download it all when I could just pull it up on youtube.''\footnote{Student W25-R041, Winter 2025 survey, improvements item.} That sentence should be read with the lecture course in the other hand, for passing-the-time is boredom's counter-move, the search for a cheaper occupation than the one that holds and empties us---and here the counter-move has been institutionalized, furnished by the course itself as a parallel resource, so that the exit from the designed world costs nothing and arrives pre-approved. The fallback answers co-occur with the friction record (installations, latency, crashes) often enough to show what occasions the exit; the arithmetic of effort shows what sustains it. A designed world in deployment does not compete with idleness. It competes with the cheapest adequate substitute for itself, and every increment of threshold-labor is an argument for the substitute.

% ============================================================================
% CORRECTION FLAG (2026-07-13, author input) — FACTUAL ERROR in this paragraph.
% Claims the async course "starves the WE" and that co-disclosedness is "less an
% experience than a demand," "the course's missing WE." THIS IS WRONG:
%   (1) the platform HAS spatial proximity-based voice chat (co-watching in-world);
%   (2) EVERY student is assigned to a procedure-paired group.
% WE-affordance and grouping both already exist. Real gap = STUDENT-SIDE ACTIVATION
% (convening the assigned group synchronously). "Watch with others" = the WE HAD by
% activators; "seemed isolating" = a non-activator who defaulted to solo;
% "more collaboration" = request to MANDATE/SCHEDULE the meet, not to supply capability.
% DO NOT COMPILE AS-IS. Rewrite pending author sign-off. See survey-analysis/05-findings.md F7.
% ============================================================================
What the world gave those who stayed is just as legible, and it answers the dissertation's second goal directly. Presence-language belongs to the app-users---4.3 mentions per hundred of their responses against 1.2 for the video-only---and one student states the two-pole account almost in the dissertation's own words: ``The virtual space feels like you are in it, while just a video doesn't. It makes you more involved and interactive with the space.''\footnote{Student W25-R044, Winter 2025 survey, use-cases item.} That is world-disclosedness in a desktop rendering: not a claim that the screen route equals the headset, but testimony that a navigable somewhere, even flat, opens as a place in a way a video does not. The other pole appears mostly as an absence with a voice. Shared-involvement talk is stable across the terms and dominantly optative---students asking for collaboration, for discussion, for company: one wants ``Having more of a community since it just seemed isolating.''\footnote{Student W25-R089, Winter 2025 survey, improvements item.} in a course that, run asynchronously, starves the WE; another names the platform's one social affordance as its chief virtue---``It is more effective that it does put you into a VR clinical environment and a cool interactive space. The main pro to the platform is you can watch with others.''\footnote{Student W26-R037, Winter 2026 survey, use-cases item.} Co-disclosedness, in this record, is less an experience than a demand: the world's multi-user lobby is the one room where the course's missing WE could live, and the students know it. And some know the gradient, too, locating the missing intensity themselves: ``This class should provide students with VR headsets. Even if there were a technology fee, it would be more beneficial for learning.''\footnote{Student W26-R014, Winter 2026 survey, platform item.} The requests run exactly along the axis the dissertation's account of the image would predict---the flatter the route, the more the student's own \textit{phantasia} must carry, and some students feel the load and ask the institution to shoulder it.

The deployment's deepest signal, though, is not a friction and not a request; it is a change in the kind of complaint. Across the three terms the crash record collapses and stays low---17, then 3, then 4---while value-skepticism ends the era threefold higher than it began: 1.3, then 0.7, then 4.2 answers per hundred, the composition caveat duly entered, the cross-term presence of the theme secure. As the platform stabilized, that is, the objection migrated from \textit{it does not work} to \textit{why does it exist}---``I think the youtube videos were enough to fully grasp the experience.''\footnote{Student W26-R033, Winter 2026 survey, platform item.}---and the divergence register gives the migration its structure. Eighteen answers praise and hollow in the same breath. Most of them, held to the diagnostic grid, resolve into the first form after all, for they name their culprit even while praising the idea: ``The actual idea of having the VR clinic is amazing, the videos and images that come with it not so much''\footnote{Student S25-R022, Spring 2025 survey, platform item.} A residue does not. ``The virtual reality environment, while it allows for more immersion, wasn't the highlight of the course, and could be amended for something else.''\footnote{Student W26-R016, Winter 2026 survey, improvements item.}---more immersion conceded, no defect named, the highlight simply not there: an emptiness without a culprit, which is the second form's exact grammar, met here as an echo in a single voice and claimed as nothing more. One student even supplies the divergence as a pair, praising the world's placement and its company in one answer and flattening it in the next: ``The VR platform is really laggy at times and is not significantly more immersive then a YouTube video''\footnote{Student W26-R037, Winter 2026 survey, platform item.} A single channel cannot detect the second form; the method owns that ceiling. But the arc from broken to pointless is the silhouette that form casts at cohort scale---everything runs, and it is still empty---and it fixes the design brief that the final section takes up. Repair answers the first form. Only a reason answers the second: a for-the-sake-of-which that makes this world, rather than its cheapest substitute, the place where the learning happens.

\section{Why Elements Bore: The Corpus Read as One}

Set the laboratory and the deployment side by side and one fact towers over the rest, a fact about instruments before it is a fact about boredom: wherever an episode of this corpus turned troubled and its channels were permitted to speak separately, the channels spoke differently. The divergence appears at four scales, and at each scale it keeps the same shape. Between instruments: on the same clinical film, the cortical profile read boring-like while the gaze read engaging-like with the corpus's one emphatic significance. Within one instrument: a participant, given a boredom question and an engagement question rather than a single scale, rated the same film at boredom six and engagement seven, and the next at seven and six. Within one student, across two answers of one survey: the platform praised for placing the student in a clinical world with company, and flattened, an item later, as not significantly more immersive than a video. Within one sentence: an idea called amazing and its execution called empty, one breath apart. Four scales, one grammar---and never the reverse grammar, never a hollow surface over a glowing depth; the divergence always stacks occupation above emptiness. That patterning is what rules out measurement noise, which has no preferred direction, and it is what convicts the bipolar instrument, which has no way to record a person occupied and empty at once. The corpus did not merely fail to fit a single engaged-to-bored axis; it refuted the axis four ways, each time by the simple device of letting two layers answer separately. If the reanalysis had produced nothing else, it would have produced this: engagement and boredom are not the two ends of one line but the two layers of one episode, and only their disagreement is diagnostic.

The second fact is older, and the corpus confirms it with almost comic thoroughness: boredom is never a property of the stimulus. The word-processor tutorial---the program's designated dullness, chosen to bore---bored seven of eight, two of them within forty seconds; the eighth, arriving depleted from the film before, met it as relief, rated it boredom one and engagement seven, and left restored, fatigue fallen from seven to one---while the interesting film, chosen to engage, bored one participant within a minute. The clinical footage lived the most divided life of all: the most presence-praised content of the deployment, the site of the laboratory's deepest divergence, and, for nearly half of every deployment term, something met only as a plain video after all. The same footage bores, engages, restores, and gets refused, according to the episode into which it arrives---the viewer's fatigue, the task's presence or absence, the hour's history. Nothing here surprises the frame, which never located boredom in things; but it is worth pausing on how completely the record cooperates. The cause-effect schema, which would assign each film a fixed valence, cannot survive its own control conditions; and once it falls, the question this section answers changes form. Not: which stimuli are boring? But: what must an encounter lack, for boredom to arrive in it?

The corpus answers with the two moments, each wearing an empirical face, and the first of them---being left empty---is the world's side: solicitation withheld. In the laboratory it looks like interest collapsing at the threshold---the offer refused within minutes, and everything after the refusal mere endurance; in the deployment it looks like the withheld object, hours of footage with nothing determinate to want from them, screens hard to read in rooms hard to see, a world that gives the eye no assignment. Being held in limbo is time's side: duration mandated, completion deferred. In the laboratory it is the while stretching past its clock in sixteen of twenty-four episodes---and stretching, note, across every stimulus category, for the strapped-in, film-mandated situation is itself a small architecture of holding; in the deployment it is the several seconds of buffer between a student's selection and the image's arrival, limbo in miniature, repeated at each threshold of intent. Boredom, on the corpus's showing, arrives exactly where the two moments meet: held before an offer that gives nothing to complete. Neither moment suffices alone, for an empty offer that can be left is merely left, and a full offer that holds is called absorption. It is the conjunction---the being-kept in the presence of the withheld---that the record tracks, in drag ratings, in sleep-fights, in the felt hour of a seventeen-minute film.

And the conjunction, in a designed world, has an address. Follow the chain: the deployment's environment produces perceptual arrival, produces the charged image, produces interest---159 mentions of the content's value, and every co-occurrence of incorporation-language in the design record landing on content rather than on action, twenty-one against zero---and then stops. Things to do: 36 mentions, not one in the company of incorporation-language; requests for interactivity: 51. The laboratory's clinical film has the same anatomy: interest ratings high, and boredom's onset by five minutes in most participants anyway, because interest without a task is a chain idling at its penultimate link. The stall's address, in the dissertation's terms, is the transition from completed cognition to completed act---from the world understood to the world acted-in---and its design lesson is blunt. A world that can only be watched has a boredom half-life, however good the watching; solicited attention that is given no act in which to complete itself will curdle by the clock; and the difference between a designed world and a film is precisely that the world could have asked for something.

Boredom's counter-move is just as design-visible, and the corpus displays it at three scales, ascending. At the scale of the episode: the searching gaze, high variance roving for anything to attend to, passing-the-time written in oculometry---and its collapse, the stilled stare, when even the search is abandoned. At the scale of the course: the fallback, forty-seven students who could have entered the world doing the effort-arithmetic and choosing its cheapest adequate substitute, an exit the course itself had furnished and pre-approved. At the scale of the institution: the turn itself, presence traded for feasibility in the move from headset to screen---the course, as a designed thing, performing the same economizing flight its students perform. Passing-the-time is not noise around this record; it is the record's shape, and each scale of it is a datum about design, since every exit taken maps a place where the world's offer lost to its overhead. What the exits could not explain, the late testimony does. As the platform stabilized---crashes collapsing after the first term---the complaint changed kind, from it does not work to why does it exist, and the register of praise-with-emptiness sorted, under the grid, into culprit-cases and a small culprit-less residue. That residue, and the arc it crowns, is the second form's silhouette at cohort scale: everything runs, and it is still empty. The method's ceiling stands---one channel, echo not detection---but the silhouette is enough to fix the design brief, because it names the one complaint that no repair addresses. The first form asks the designer for fixes. The second asks the designer for a reason.

Two threads complete the account, closing it at its proper depth, and the first of them is the aftermath: boring episodes deplete--- fatigue at or above its start in six of eight after the tutorial, twice climbing from a three to a seven and an eight; a study reordered around the exhaustion its own control induced; and the psychology of attention supplying the mechanism's name, attention residue---the prior content interfering with the present attending, most where the present is least engaging.\footnote{Mark, \textit{Attention Span}, 97.} The lecture course, which reads boredom as disclosive, never theorizes what boredom costs, and here the corpus sends something back: whatever boredom shows, it also spends, and a curriculum that ignores the spending will pay it in the next module. The second thread is the gradient's correction, for the deployment's students testify to the image's economy exactly as the dissertation would have it---presence attested on the desktop route, the flatter route felt as thinner, headsets requested by name---and yet the skepticism rose as the platform ran smoothest, because intensity of image, alone, justified nothing. The route sets the load; only the task sets the dwelling; and a world upgraded to perfect vividness, with nothing to do and no one to do it with, would merely be bored in at higher resolution. The scope of all this is owned before the close: one course, one platform, laboratory cohorts of three, eight, and twelve---structural claims drawn from convergent cases, never population estimates. That is why the account must end past the empirical record, at the formulation the whole corpus has been circling. ``Boredom is the entrancement of the temporal horizon, an entrancement which lets the moment of vision belonging to temporality vanish.''\footnote{Heidegger, \textit{Fundamental Concepts}, 153 / GA 29/30: 230.} What the bored episodes of this record share---across their different faces: the searching eye, the stretched while, the occupied-and-empty hour, the flight to the cheaper substitute---is time entrancing its holder, and the vanishing of the moment in which something could be seized and done. The lecture course adds, in its deepest and most cautious claim, that ``the third form is the condition of the possibility of the first and thereby also of the second''\footnote{Heidegger, \textit{Fundamental Concepts}, 156--57 / GA 29/30: 235.}---the shallower boredoms nested in the deeper possibility, an ordering this Part carries only as structure, under its stated analogy. Structure is enough. A designed world cannot manufacture the moment of vision; that belongs to the student. What it can do---what the whole record says it must do---is stop letting the moment vanish by default: give the held attention an act, give the act a companion, and give the hour a reason. How, is the last section's business.

\section{Design for Dasein in the Classroom: Toward Better Virtual Learning Environments}

A diagnosis that stopped at diagnosis would betray the corpus that produced it, for the same record that shows where the chain stalls shows, by the same light, where a designer must go to unstall it, and the frame has held from the beginning that these are one knowledge and not two. The design lens brought to this Part was built on exactly that identity: to design a virtual environment is to design the conditions under which a participant's actualization of desire runs or stalls---the conditions, never the outcome, since a world can occasion an act but never cause one---and designing under that description is less an engineering afterthought than philosophy continued by other means, ``doing philosophy with the hands''\footnote{Thomas Wendt, \textit{Design for Dasein: Understanding the Design of Experiences} (self-published, 2015), 6.} in the design theorist's own phrase. What the empirical Part adds to the stance is an address book. Where the earlier sections of the dissertation could say only that a designed world solicits, holds, empties, or completes its participants in general, the corpus now says where this world did each of those things to these students, in what numbers, and with what testimony; and a design brief written from such a record inherits its evidence the way a treatment inherits an examination. The five levers that follow are therefore not a wish list but a ranked reading of the record---ranked by yield, cheapest and largest first, because the deployment's own history teaches that feasibility is not the enemy of the designed world but its zeroth condition.

The first lever is the threshold, and it outranks everything inside the world. Among the deployment's largest friction families are those of the approach rather than the interior---installations, hardware, a sign-in that would not admit its student, and crashes that struck at the threshold as often as within---and together they shut a quarter of the class out outright, taxed the entry of the rest, and fed the flight of forty-seven students who could have entered and chose the cheaper substitute instead. Remedy follows diagnosis: deliver the world at zero setup---browser-grade access, pre-installed laboratory machines, a build that does not care what laptop a student owns. Every increment of entry labor is an argument for the fallback, and no room, however well designed, teaches the student who never reached it. The second lever is the task, and it answers the record's deepest complaint. The world affords interest massively and action not at all---one hundred fifty-nine mentions of the content's value against thirty-six mentions of things to do, none of the latter in the company of a single incorporation-positive word, and fifty-one open requests for interactivity. So furnished, the chain runs to the threshold of completion and idles there, in a world that can only be watched. Give the world a task: needs-finding done \textit{inside} the environment, routes that must be chosen, checkpoints that must be reached, assignments whose deliverable exists only in the world. A task opens the transition from cognition to act that the record shows foreclosed. And a task, alone among the levers, answers the late question that no repair could touch, because a task is a reason---the for-the-sake-of-which that makes this world, rather than its cheapest substitute, the place where the learning happens. The third lever supplies the object of the looking. Students facing hours of footage asked, in so many words, what they were supposed to be looking for; they asked for captions and transcripts; they reported screens hard to read and procedures hard to parse. Give the looking its \textit{orekton}: viewing briefs that say what to seek and why it matters, segmented procedures with named phases, legible text on legible screens---small furnishings, but they are the difference between an offer and a blank. The fourth lever is the WE. Shared involvement appears in this record almost entirely as a want---isolation lamented, collaboration requested, watching-with-others named the platform's chief virtue by a student who otherwise flattened it. An asynchronous course starves precisely the pole of the union that the environment's multi-user rooms were built to feed. Schedule the WE: make watching-together the default first encounter, put group tasks in the world, let the course's one shared place actually be shared. And the fifth lever is the aftermath. Boredom depletes---the laboratory record shows fatigue leaving the boring film at or above its starting level in six of eight participants, and a study reordering itself around the exhaustion its own control induced. Sequence is therefore a design material: nothing hostile first, recovery after demand, module lengths that respect how early the drag begins on unfriendly material. None of these five requires a headset, a budget, or a discovery; four of the five were named, in one voice or another, by the students themselves, and the fifth was written by the laboratory record. What the frame adds is the order and the address: which stage of the chain each lever serves, and why the second, the task, is the one that no amount of polish elsewhere can replace.

The brief needs a motto, and the record has already written it: design against the question \textit{why not just YouTube?} Every element of a designed learning world should survive that interrogation---should do something the cheapest adequate substitute cannot, whether the something is a task, a place, or a company. The deployment shows, with unusual clarity, that students audit their worlds like economists, and that a world which cannot answer the audit will be fled at the first friction and doubted even when it runs. The motto has a mirror, and honesty requires facing it. The same knowledge this Part has assembled---where attention stalls, what occupies its surface while its depth empties, when a person is most susceptible to a well-aimed solicitation---is held, at industrial scale and to opposite purpose, by the attention economy, whose instruments do not repair the stalled chain but harvest it: ``The algorithms incorporate this information---who you are, how you feel, what you do when and where---and then use it to capture your attention.''\footnote{Mark, \textit{Attention Span}, 155.} The notification engineered for the bored student and the task designed for the same student are twins with opposite vocations, and the difference between them is not the knowledge but the norm. The design theory this Part leans on states the norm plainly---advocate preferable worlds; refuse the dark patterns\footnote{Wendt, \textit{Design for Dasein}, on design's advocacy of preferable worlds against dark patterns [****** page unverified at assembly].}---and the frame sharpens it: capture bypasses the participant's own efficient and final causes, while a task addresses them, and only the second is teaching. It is in that spirit, and only as a gesture toward future work, that one idea from the author's own notes belongs here: the same ambient instrumentation that the capture economy aims at engagement could run care-side---eye tracking and data logging in the background of a learning world, noticing when a student's focus behavior degrades and telling a teacher rather than an advertiser. Whether such a tool should exist is a question for the norm just stated, and its design is no part of this dissertation; the gesture marks only that the instrumentation's vocation is a choice.

A second gesture closes the empirical program's account with the frame's return-gift to measurement. If this corpus teaches one methodological lesson, it is that boredom and engagement must be listened for on two channels at once, and that the instrument to build is not a better bipolar scale but a three-state reading---restless search, occupied-but-hollow, withdrawn---taken from a surface channel and a depth channel whose \textit{disagreement} is the diagnostic. Felt duration enters as a depth probe this record has already validated, and sleep-exit belongs in any such instrument as its own terminal category rather than as boredom's maximum. The proposal's details belong to a future study; its principle belongs here, because it is the empirical face of the Part's whole argument: one axis was never enough.

What remains is the ledger of limits, and then the close. The category caution has governed every page and governs this one: nothing in a seventeen-minute film or a ten-week course discloses beings as a whole, and the profound form of boredom entered this Part only as a structure lending its grammar, never as a claim about any student's metaphysical hour. The body's channel remains untested---a desktop deployment cannot speak to locomotion, to room-scale naturalization, to what the headset's route would have asked of these students' kinesthetic learning---so the levers claim the routes on which they were measured and no others. The cohorts were small, the course was one course, and every structural claim stands on convergence rather than count. Within those bounds, the Part has done what it promised at its door. Boredom has been explained---not as a property of dull stimuli but as the meeting of a withheld offer and a held while, legible in gaze and cortex and felt time and flight; engagement has been assessed---and found to be not the opposite of boredom but its neighboring layer, occupied surface over a depth that must be listened for separately; and design has been instructed, in five levers and a motto, by the same frame that did the explaining. The deepest thing the record showed was an entrancement: time closing over its holder, the moment of vision withheld, in a laboratory chair and in a virtual clinic alike. No designer can manufacture that moment; it belongs to the student, as the act belongs to the agent, and a frame that began with the motion of the soul would forget itself if it ended by promising machinery for what is ownmost. But a designed world decides, every hour it runs, whether the moment's conditions are furnished or foreclosed---whether held attention is given an act, the act a companion, and the hour a reason---and that decision, made well or badly, is rhetoric's oldest work carried into its newest rooms: the arrangement of a world so that a soul, moved, may complete its motion. The chain this dissertation began with ends where it always ends, in an act; the dissertation ends by handing the act to the designer.

\end{document}